% =====================================================================
\chapter{Θεωρητικό Υπόβαθρο}
\label{ch:background}
% =====================================================================

Το κεφάλαιο αυτό παρουσιάζει τις έννοιες στις οποίες στηρίζεται η εργασία: την
αρχιτεκτονική του διαύλου \en{USB} και τους περιγραφείς του (§\ref{sec:usb-arch}),
τη διαδικασία απαρίθμησης (§\ref{sec:enumeration}), την κλάση \en{HID}
(§\ref{sec:hid}), την κλάση μαζικής αποθήκευσης και τους πίνακες διαμερισμάτων
(§\ref{sec:msc}), και τέλος τους μηχανισμούς του πυρήνα \en{Linux} που
αξιοποιεί το \Probolos{} (§\ref{sec:linux-usb}).

\section{Αρχιτεκτονική USB}
\label{sec:usb-arch}

Το \en{USB} είναι ένας ασύμμετρος δίαυλος τύπου \en{host--device}: ο
κεντρικός υπολογιστής (\en{host}) ελέγχει πλήρως την επικοινωνία και οι συσκευές
απαντούν μόνο όταν ερωτώνται~\cite{usb20spec}. Η τοπολογία είναι δενδρική: κάθε
ελεγκτής κεντρικού υπολογιστή (\en{host controller}) εκθέτει έναν
\emph{ριζικό διανομέα} (\en{root hub}), κάτω από τον οποίο συνδέονται συσκευές
ή επιπλέον διανομείς, έως επτά επίπεδα βάθος.

\subsection{Ιεραρχία περιγραφέων}

Κάθε συσκευή \en{USB} περιγράφει τον εαυτό της με μια ιεραρχία δομών δεδομένων
που ονομάζονται περιγραφείς (Σχήμα~\ref{fig:descriptors}):

\begin{itemize}
  \item \textbf{Περιγραφέας συσκευής} (\en{device descriptor}, 18 \en{bytes}):
        αναγνωριστικά κατασκευαστή και προϊόντος (\code{idVendor},
        \code{idProduct}), έκδοση (\code{bcdDevice}), έκδοση προδιαγραφής
        (\code{bcdUSB}), κλάση συσκευής και αριθμό διαμορφώσεων.
  \item \textbf{Περιγραφέας διαμόρφωσης} (\en{configuration descriptor}):
        αριθμός διεπαφών, χαρακτηριστικά τροφοδοσίας και μέγιστο ρεύμα
        (\code{bMaxPower}, σε μονάδες των 2\,\en{mA} για \en{USB} 2 και 8\,\en{mA}
        για \en{SuperSpeed}). Το πεδίο \code{wTotalLength} δηλώνει το συνολικό
        μήκος της αλυσίδας περιγραφέων της διαμόρφωσης.
  \item \textbf{Περιγραφέας διεπαφής} (\en{interface descriptor}): η τριάδα
        κλάσης/υποκλάσης/πρωτοκόλλου (\code{bInterfaceClass},
        \code{bInterfaceSubClass}, \code{bInterfaceProtocol}) που καθορίζει ποιος
        οδηγός θα προσδεθεί.
  \item \textbf{Περιγραφείς τελικών σημείων} (\en{endpoint descriptors}): τα
        κανάλια επικοινωνίας και ο τύπος μεταφοράς τους.
  \item \textbf{Περιγραφείς συμβολοσειρών} (\en{string descriptors}): κείμενα
        \en{UTF-16LE} για κατασκευαστή, προϊόν και σειριακό αριθμό, που
        \emph{επιλέγονται ελεύθερα από τη συσκευή}.
\end{itemize}

\begin{figure}[htbp]
\centering
\begin{tikzpicture}[
  box/.style={draw, rounded corners=2pt, minimum width=3.2cm, minimum height=0.85cm,
              align=center, font=\small, fill=#1},
  box/.default=white,
  >={Stealth[length=2.2mm]}]
  \node[box=univblue!12] (dev) {Περιγραφέας συσκευής\\\footnotesize\en{VID, PID, bcdUSB}};
  \node[box=univblue!6, below=0.8cm of dev] (cfg) {Διαμόρφωση 1\\\footnotesize\en{bMaxPower}};
  \node[box=okgreen!10, below left=0.8cm and 0.2cm of cfg] (if0) {Διεπαφή 0\\\footnotesize\en{Class 08: Mass Storage}};
  \node[box=critical!10, below right=0.8cm and 0.2cm of cfg] (if1) {Διεπαφή 1\\\footnotesize\en{Class 03: HID / Keyboard}};
  \node[box, below=0.7cm of if0, minimum width=1.5cm] (ep0) {\footnotesize\en{EP 1 IN/OUT}};
  \node[box, below=0.7cm of if1, minimum width=1.5cm] (ep1) {\footnotesize\en{EP 2 IN}};
  \node[box=gray!10, right=1.2cm of dev] (str) {Συμβολοσειρές\\\footnotesize\en{"SanDisk Cruzer"}};
  \draw[->] (dev) -- (cfg);
  \draw[->] (cfg) -- (if0);
  \draw[->] (cfg) -- (if1);
  \draw[->] (if0) -- (ep0);
  \draw[->] (if1) -- (ep1);
  \draw[->, dashed] (dev) -- (str);
\end{tikzpicture}
\caption{Ιεραρχία περιγραφέων μιας σύνθετης συσκευής. Η συγκεκριμένη δηλώνει
ταυτόχρονα διεπαφή μαζικής αποθήκευσης και πληκτρολογίου — τυπική μορφή
συσκευής \en{BadUSB}.}
\label{fig:descriptors}
\end{figure}

\subsection{Κλάσεις συσκευών}

Η κλάση μιας διεπαφής καθορίζει τη λειτουργία της. Ο Πίνακας~\ref{tab:usb-classes}
παραθέτει τις κλάσεις που είναι κρίσιμες για την παρούσα εργασία. Μια συσκευή με
περισσότερες από μία διεπαφές διαφορετικών κλάσεων ονομάζεται \emph{σύνθετη}
(\en{composite}). Οι σύνθετες συσκευές είναι νόμιμες και συνηθισμένες (π.χ.
\en{webcam} με μικρόφωνο), αλλά ορισμένοι συνδυασμοί δεν έχουν νόμιμη χρήση.

\begin{table}[htbp]
\centering
\caption{Επιλεγμένες κλάσεις διεπαφών \en{USB} και η σημασία τους για την ασφάλεια.}
\label{tab:usb-classes}
\begin{tabular}{@{}clp{7.5cm}@{}}
\toprule
\textbf{Κωδ.} & \textbf{Κλάση} & \textbf{Σημασία για την ασφάλεια}\\
\midrule
\code{0x01} & \en{Audio}           & χαμηλός κίνδυνος· πιθανό κανάλι υποκλοπής (μικρόφωνο)\\
\code{0x02} & \en{CDC}             & επικοινωνία/δίκτυο· μπορεί να ανακατευθύνει κίνηση\\
\code{0x03} & \en{HID}             & πληκτρολόγιο/ποντίκι· \textbf{μπορεί να πληκτρολογεί}\\
\code{0x08} & \en{Mass Storage}    & αποθήκευση· φορέας κακόβουλων αρχείων/συστημάτων αρχείων\\
\code{0x09} & \en{Hub}             & διανομέας· επεκτείνει την τοπολογία\\
\code{0x0E} & \en{Video}           & κάμερα· κανάλι υποκλοπής\\
\code{0xE0} & \en{Wireless}        & ραδιοσυσκευές (π.χ. \en{Bluetooth})· γέφυρα σε ασύρματες συσκευές\\
\code{0xFF} & \en{Vendor-specific} & άγνωστη λειτουργία· δεν επιδέχεται σημασιολογικό έλεγχο\\
\bottomrule
\end{tabular}
\end{table}

\section{Η διαδικασία απαρίθμησης}
\label{sec:enumeration}

Όταν μια συσκευή συνδέεται, ο διανομέας ανιχνεύει τη μεταβολή και ο ελεγκτής
ξεκινά την \emph{απαρίθμηση} (\en{enumeration})~\cite{usb20spec}:

\begin{enumerate}
  \item Επαναφορά της θύρας (\en{reset}) και ανίχνευση ταχύτητας.
  \item Ανάγνωση των πρώτων 8 \en{bytes} του περιγραφέα συσκευής, για το
        \code{bMaxPacketSize0}.
  \item Ανάθεση διεύθυνσης (\code{SET\_ADDRESS}).
  \item Ανάγνωση του πλήρους περιγραφέα συσκευής και όλων των διαμορφώσεων.
  \item Ανάγνωση των περιγραφέων συμβολοσειρών.
  \item Επιλογή διαμόρφωσης (\code{SET\_CONFIGURATION}) και πρόσδεση οδηγών
        στις διεπαφές.
\end{enumerate}

Κρίσιμη παρατήρηση: τα βήματα 1–5 εκτελούνται \emph{πάντα} από τον πυρήνα, ενώ
το βήμα 6 εξαρτάται από την εξουσιοδότηση της συσκευής (§\ref{sec:authorization}).
Συνεπώς το \Probolos{} μπορεί να διαβάσει όλους τους περιγραφείς μιας συσκευής
χωρίς αυτή να έχει αποκτήσει οδηγό. Αντίστροφα, οποιοδήποτε σφάλμα στον κώδικα
απαρίθμησης του ίδιου του πυρήνα είναι εκτός του πεδίου προστασίας του εργαλείου.

\section{Συσκευές Διεπαφής Ανθρώπου (HID)}
\label{sec:hid}

Η κλάση \en{HID}~\cite{usbhid111} καλύπτει πληκτρολόγια, ποντίκια, χειριστήρια
και πλήθος άλλων συσκευών. Μια συσκευή \en{HID} περιγράφει τη μορφή των
αναφορών της μέσω ενός \emph{περιγραφέα αναφορών} (\en{report descriptor}). Για
λόγους εκκίνησης (\en{BIOS}), ορίζεται επίσης το «πρωτόκολλο εκκίνησης»
(\en{boot protocol}), όπου η υποκλάση \code{1} και το πρωτόκολλο \code{1}
δηλώνουν πληκτρολόγιο και το \code{2} ποντίκι.

Από τη σκοπιά της ασφάλειας, η κλάση \en{HID} είναι ιδιαίτερα επικίνδυνη διότι:
(α) δεν απαιτεί πρόσθετο οδηγό σε κανένα λειτουργικό σύστημα· (β) ό,τι
«πληκτρολογεί» η συσκευή γίνεται δεκτό ως είσοδος του χρήστη, με τα δικαιώματα
της ενεργής συνεδρίας· και (γ) ο ρυθμός πληκτρολόγησης μπορεί να ξεπερνά τους
1000 χαρακτήρες ανά δευτερόλεπτο. Η επίθεση \emph{έγχυσης πληκτρολογήσεων}
(\en{keystroke injection}) αξιοποιεί ακριβώς αυτές τις ιδιότητες.

\section{Μαζική αποθήκευση και πίνακες διαμερισμάτων}
\label{sec:msc}

Οι συσκευές μαζικής αποθήκευσης (\en{MSC}) υλοποιούν συνήθως το πρωτόκολλο
\en{Bulk-Only Transport} με εντολές \en{SCSI}~\cite{usbmsc}. Στο \en{Linux}
εμφανίζονται ως συσκευές μπλοκ (\code{/dev/sdX}). Το περιεχόμενό τους
οργανώνεται με πίνακα διαμερισμάτων:

\begin{itemize}
  \item \textbf{\en{MBR}}: στον τομέα 0, με τέσσερις εγγραφές των 16 \en{bytes}
        (τύπος, αρχικό \en{LBA}, μήκος) και την υπογραφή \code{0x55AA}.
  \item \textbf{\en{GPT}}~\cite{uefi210}: προστατευτικό \en{MBR} στον τομέα 0,
        κεφαλίδα στο \en{LBA}~1 και εγγραφές διαμερισμάτων συνήθως από το
        \en{LBA}~2, με \en{GUID} τύπου και χαρακτηριστικά για κάθε διαμέρισμα.
\end{itemize}

Ένα διαμέρισμα συστήματος \en{EFI} (\en{ESP}) ή ένα «κρυφό» διαμέρισμα σε μια
κάρτα μνήμης δεν έχει συνήθως λόγο ύπαρξης και μπορεί να υποδεικνύει
προσπάθεια εκκίνησης κακόβουλου κώδικα ή απόκρυψης δεδομένων.

\section{Το υποσύστημα USB του Linux}
\label{sec:linux-usb}

\subsection{Το \en{sysfs} και η εξουσιοδότηση συσκευών}
\label{sec:authorization}

Ο πυρήνας \en{Linux} εκθέτει κάθε συσκευή \en{USB} ως κατάλογο στο
\pth{/sys/bus/usb/devices/}. Ήδη από την έκδοση 2.6.24 υποστηρίζεται η
\emph{εξουσιοδότηση} συσκευών, αρχικά για το \en{Wireless USB}~\cite{kernel-usb-authorization}:

\begin{itemize}
  \item \code{authorized} (ανά συσκευή): 1 αν η συσκευή επιτρέπεται να
        διαμορφωθεί και να αποκτήσει οδηγούς, 0 αν όχι. Η εγγραφή 0 σε μια
        ενεργή συσκευή αποσυνδέει τους οδηγούς της.
  \item \code{authorized\_default} (ανά ριζικό διανομέα): η τιμή
        \code{authorized} που λαμβάνουν οι \emph{νέες} συσκευές. Η τιμή 2
        σημαίνει «μόνο εσωτερικές θύρες».
  \item \code{authorized} ανά διεπαφή: επιτρέπει την απενεργοποίηση μεμονωμένων
        διεπαφών μιας συσκευής.
\end{itemize}

\begin{lstlisting}[language=shell,caption={Χειροκίνητο «κλείσιμο» και άνοιγμα της πύλης μέσω \en{sysfs}.},label=lst:sysfs-manual]
# Close: new devices arrive unauthorized
for hub in /sys/bus/usb/devices/usb*/authorized_default; do
    echo 0 | sudo tee "$hub"
done
# Authorize a single device by its sysfs name (e.g. 1-2)
echo 1 | sudo tee /sys/bus/usb/devices/1-2/authorized
# Reopen
for hub in /sys/bus/usb/devices/usb*/authorized_default; do
    echo 1 | sudo tee "$hub"
done
\end{lstlisting}

Η τιμή \code{authorized\_default} είναι \emph{καθολική ανά διανομέα} και
\emph{διατηρείται} όσο λειτουργεί το σύστημα. Αν ένα πρόγραμμα την ορίσει σε 0
και τερματιστεί απρόσμενα, κάθε μελλοντική συσκευή — συμπεριλαμβανομένου ενός
πληκτρολογίου \en{USB} — θα παραμείνει νεκρή. Ο κίνδυνος αυτός «κλειδώματος
εκτός» (\en{lockout}) διαμορφώνει σε μεγάλο βαθμό τη σχεδίαση του \Probolos{}
(§\ref{sec:design-failsafe}).

\subsection{Το \en{udev}}

Ο δαίμονας \en{udev} λαμβάνει από τον πυρήνα γεγονότα \en{uevent} (\code{add},
\code{change}, \code{remove}, \code{bind}) για κάθε μεταβολή συσκευής, μαζί με
έναν αύξοντα αριθμό \code{SEQNUM}. Το \Probolos{} παρακολουθεί αυτά τα γεγονότα
μέσω της βιβλιοθήκης \en{pyudev}, ενώ κανόνες \en{udev} μπορούν να ορίσουν
ιδιότητες όπως \code{UDISKS\_IGNORE=1} για να αποτρέψουν την αυτόματη προσάρτηση.

\subsection{Το υποσύστημα εισόδου \en{evdev}}

Οι συσκευές εισόδου εκτίθενται ως κόμβοι \code{/dev/input/eventN}~\cite{kernel-input}.
Κάθε γεγονός (\code{struct input\_event}) φέρει χρονοσφραγίδα πυρήνα, τύπο,
κωδικό και τιμή. Η κλήση \code{ioctl(fd, EVIOCGRAB, 1)} δίνει σε μια διεργασία
\emph{αποκλειστική} πρόσβαση στα γεγονότα της συσκευής: κανείς άλλος αναγνώστης,
ούτε ο εξυπηρετητής γραφικών, δεν τα λαμβάνει. Αυτός ο μηχανισμός αποτελεί τη
βάση της καραντίνας συμπεριφοράς (§\ref{sec:design-stage3}).

\subsection{Εξομοίωση συσκευών: \en{dummy\_hcd} και \en{raw\_gadget}}

Ο πυρήνας περιλαμβάνει έναν εικονικό ελεγκτή (\en{dummy\_hcd}) που συνδέει τη
στοίβα «συσκευής» (\en{gadget}) με τη στοίβα «κεντρικού υπολογιστή» του ίδιου
μηχανήματος, και τη διεπαφή \en{raw\_gadget}~\cite{kernel-raw-gadget} που
επιτρέπει την υλοποίηση συσκευών \en{USB} από χώρο χρήστη. Ο συνδυασμός τους
επιτρέπει τη δημιουργία αυθαίρετων — και κακόβουλων — συσκευών \en{USB} χωρίς
υλικό, κάτι που αξιοποιείται στην αξιολόγηση (§\ref{sec:eval-testbed}).

\section{Αρχές σχεδίασης ασφαλών συστημάτων}
\label{sec:principles}

Η σχεδίαση στηρίζεται στις κλασικές αρχές των \en{Saltzer} και
\en{Schroeder}~\cite{saltzer1975protection}:

\begin{itemize}
  \item \textbf{Προεπιλογές ασφαλείς στην αποτυχία} (\en{fail-safe defaults}):
        η πρόσβαση στηρίζεται σε ρητή άδεια, όχι σε αποκλεισμό.
  \item \textbf{Ελάχιστο προνόμιο} (\en{least privilege}): κάθε τμήμα έχει μόνο
        τα δικαιώματα που χρειάζεται.
  \item \textbf{Οικονομία μηχανισμού} (\en{economy of mechanism}): η έμπιστη
        βάση πρέπει να είναι όσο το δυνατόν μικρότερη.
  \item \textbf{Πλήρης διαμεσολάβηση} (\en{complete mediation}): κάθε πρόσβαση
        ελέγχεται.
  \item \textbf{Ψυχολογική αποδοχή} (\en{psychological acceptability}): ο
        μηχανισμός πρέπει να είναι εύχρηστος, αλλιώς θα παρακαμφθεί.
\end{itemize}

Ο διαχωρισμός προνομίων, όπως τον εφάρμοσαν οι \en{Provos et~al.} στο
\en{OpenSSH}~\cite{provos2003privsep}, χωρίζει ένα πρόγραμμα σε μικρό
προνομιούχο τμήμα και μεγαλύτερο μη προνομιούχο τμήμα που επεξεργάζεται μη
έμπιστα δεδομένα· η ίδια ιδέα εφαρμόζεται στο \Probolos{} (§\ref{sec:design-privsep}).
